\section{指数与对数：相同相对比例的变化}
\label{sec:02-Calculus/01-Functions-and-Precalculus/04}

两条训练损失曲线画在同一张图上。跑到后半程，两条都紧贴横轴，看上去都停住了，于是有人建议把训练砍掉一半。

把纵轴换成对数刻度再看：一条仍在稳稳往下掉，斜率几乎是常数；另一条早就平了。同一份数字，只换了纵轴的刻度，结论相反。

换刻度为什么能改变判断？因为线性纵轴回答的是``减少了多少''，而从 \(0.010\) 掉到 \(0.001\) 在这个尺度上和从 \(0.100\) 掉到 \(0.091\) 一样不起眼——尽管前者是掉到十分之一，后者只掉了不到一成。对数纵轴换了一个问题：减少了几倍。本节要处理的整族函数，就是专门描述``每一步乘以同一个数''这种变化的。

\subsection{固定增量与固定倍数是两种世界}

每步加 \(5\)，得到等差；每步乘 \(1.05\)，新增量随当前规模一起长大，得到等比。离散时刻下后者写成 \(y_n=y_0a^n\)，\(a>1\) 增长，\(0<a<1\) 衰减。

要让时间连续取值，就得把 \(a^x\) 从整数指数扩到实数指数，唯一的约束是别把加法搬成乘法这条性质弄丢：

\[
a^{x+y}=a^x a^y .
\]

补全之后的指数函数对全体实数有定义，取值恒正，穿过 \((0,1)\)，以横轴为渐近线。恒正这一点后面会不断被用到——它意味着指数函数没有零点，也意味着 \(e^{\text{任何东西}}\) 可以安全地放进分母。

指数增长有多不讲道理，画出来最清楚。

\begin{center}
\begin{tikzpicture}[font=\small,>=stealth]
\begin{scope}[x=0.55cm,y=1cm]
  \draw[->] (-0.2,0) -- (7.7,0) node[right,font=\footnotesize] {\(x\)};
  \draw[->] (0,-0.2) -- (0,4.5) node[above,font=\footnotesize] {\(y\)};
  \draw[blue!70!black,thick] plot[domain=0:7,samples=60] (\x,{\x*\x/32});
  \draw[red!70!black,thick] plot[domain=0:7,samples=60] (\x,{exp(0.6931*\x)/32});
  \node[font=\footnotesize,blue!70!black] at (5.9,0.85) {\(x^2\)};
  \node[font=\footnotesize,red!70!black] at (5.0,3.4) {\(2^x\)};
  \node[font=\footnotesize,black!65] at (3.6,-0.55) {线性纵轴};
\end{scope}
\begin{scope}[xshift=6.4cm,x=0.55cm,y=0.5cm]
  \draw[->] (-0.2,0) -- (7.7,0) node[right,font=\footnotesize] {\(x\)};
  \draw[->] (0,-2.8) -- (0,8.2) node[above,font=\footnotesize] {\(\log_2 y\)};
  \draw[red!70!black,thick] (0,0) -- (7,7);
  \draw[blue!70!black,thick] plot[domain=0.45:7,samples=60] (\x,{2*ln(\x)/ln(2)});
  \node[font=\footnotesize,red!70!black] at (5.0,7.2) {\(2^x\)};
  \node[font=\footnotesize,blue!70!black] at (6.3,4.3) {\(x^2\)};
  \node[font=\footnotesize,black!65] at (3.6,-1.1) {对数纵轴};
\end{scope}
\end{tikzpicture}

\emph{左图里两条曲线在 \(x=4\) 之前难分伯仲，过了这一点指数彻底甩开幂函数，再往后幂函数被压成一条贴地的线，什么也读不出来。右图把纵轴换成 \(\log_2\)：指数被拉成一条直线，斜率就是每步乘的倍数；幂函数则弯着上升。看到直线就该想到固定倍数，看到弯曲就该问倍数在怎样变化。}
\end{center}

这张图解释了工程里两个常见动作。绘图时给纵轴加对数刻度，是为了让``每轮乘以多少''这件事变成可以用尺子量的斜率；模型规模和算力横跨好几个数量级时，两个轴都取对数，幂律关系同样会被拉直成直线。反过来说，看到一条在双对数坐标下笔直的线，你得到的是``大致服从幂律''这条线索，而不是证明。

\subsection{\texorpdfstring{自然底数 \(e\) 让速率写起来最省事}{自然底数 e 让速率写起来最省事}}

底数可以随便选，因为 \(a^t=e^{t\ln a}\) 把任何底数换算成了 \(e\)。选 \(e\) 的理由只有一条，也是决定性的一条：

\[
\frac{d}{dx}e^x=e^x .
\]

指数函数的瞬时变化率等于它自己。换成人话就是``增长速度正比于当前规模''，而这正是复利、放射衰变、种群早期扩张、误差按固定比例缩小这一大类过程共同的机制。写成微分方程是 \(y'=ky\)，解出来就是

\[
y(t)=y_0e^{kt},
\]

\(k>0\) 增长，\(k<0\) 衰减。常数 \(e\approx2.71828\) 本身来自复利频率不断加密时 \((1+1/n)^n\) 的极限，这个极限存在而不发散，是后面极限一章会正式交代的事；此处只需知道它给了指数函数一个最干净的求导公式。

\subsection{对数问的是``要乘多少次''}

对数是指数的反函数：\(a^x=y\) 与 \(x=\log_a y\) 是同一句话的两种写法。上一节说过反函数交换输入输出，所以对数的定义域是 \((0,\infty)\)，值域是全体实数——负数和零没有对数，不是禁忌，是指数函数根本没到过那里。它们的图像也就是关于 \(y=x\) 的一对镜像。

对数定律不需要单独记，它们只是指数定律翻译过来的：

\[
\log_a(xy)=\log_a x+\log_a y,
\qquad
\log_a(x^r)=r\log_a x\quad(x>0).
\]

对数把乘法降成加法，把幂降成乘法。它不拆加法，\(\log(x+y)\) 与 \(\log x+\log y\) 没有任何关系，这是最常被写错的一条。换底公式 \(\log_a x=\ln x/\ln a\) 说明不同底数只差一个常数因子，即只改变纵向尺度；因此底数的选择等于单位的选择，信息论用底 \(2\) 得到比特（bit），用 \(e\) 得到奈特（nat），换算就是乘一个 \(\ln 2\)（见\hyperref[ch:054]{``信息与熵：怎样把``不确定''变成可计算的量''}）。

``降一级''这件事在数值计算里救命。一千个 \(10^{-8}\) 量级的概率连乘，双精度浮点会直接给出 \(0\)，后续所有比较全部失效；取对数之后连乘变成求和，量级落在几千这个安全范围内，比较大小的结论完全不变，因为对数是严格递增的。分类模型的损失写成对数似然而不是似然，采样算法里到处是对数概率，理由都在这里。真要把它们加起来而不是乘起来时，还需要一点额外技巧防止指数溢出，那属于\hyperref[ch:067]{``机器学习中的稳定计算''}一章的内容。

\subsection{倍增时间与初始量无关}

在 \(y(t)=y_0e^{kt}\) 里问``翻倍要多久''，解 \(e^{kT}=2\) 得 \(T=\ln 2/|k|\)；问``减半要多久''，答案形状完全一样。初始量 \(y_0\) 被约掉了，这是``相对速率恒定''的直接后果，也是半衰期这个概念能脱离样本量单独报出的原因。

同一个公式反过来用也很顺手。学习率按 \(\eta_t=\eta_0e^{-\lambda t}\) 衰减时，你真正关心的往往不是 \(\lambda\) 而是``几步之后降一半''，两者由 \(T=\ln 2/\lambda\) 互相换算，后者更容易和训练总步数对齐着调。已知细菌每 \(3\) 小时翻一番，问长到十倍要多久，同样是把未知的指数用对数拉回一次方程：\(2^{t/3}=10\) 给出 \(t=3\ln 10/\ln 2\approx 9.97\) 小时。

\subsection{固定倍数是一个假设，不是定律}

指数模型只在``相对增长率不变''成立时才成立，而这个前提在现实里几乎总有有效期。资源会耗尽，市场会饱和，政策会改变，测量口径会调整。短期拟合得好不等于可以外推——把一条早期指数曲线延长几个数量级，得到的往往是荒谬的数字，而错误不在数学里，在那句没被检查的假设里。

诊断办法就是本节开头那张图：把数据画到对数纵轴上，如果直线开始向下弯，相对速率已经在衰减，该换模型了。

指数刻画的是一路朝一个方向走的变化，越走越远。可现实里还有另一类同样常见的模式——转一圈回到原处，然后再来一遍。下一节从单位圆开始，把``重复''本身变成函数。
